\subsection{Enumeration and Boolean Formulas for Multigraphs}
\label{app:multi-finite}
\textbf{Section~\ref{sec:exact}} gives the degree and adjacency constraints used below.
For each case, the verification enumerates the possible integer data
or checks a refutation of the Boolean formula describing those
constraints.

\subsubsection{Maximum Degrees Five and Seven}\label{app:multi-low-degree}

For the 21 configurations in \textbf{Figure~\ref{fig:finite-five-spokes}},
\path{d5_parallel_spoke_bound.py} fixes the six spokes and the edges
between their nonroot endpoints. For one vertex outside the root
neighborhoods it enumerates the vectors $z$ of marked multiplicities
and $u$ of required host incidences as in
\textbf{Subsection~\ref{subsec:exact-small-degree}}. Starting with
$V(0)=0$, the recurrence bounds the additional marked edges at each
remaining capacity vector. Adding the fixed marked edges gives the
21 bounds printed in the figure, whose maximum is 27.

At $D=7$, Section~3 of \path{delta7_parallel_reduction.md} gives the
parameters for the recurrence implemented in
\path{d7_bipartite_ore_coupled.py}. For $q=9,10$, the resulting bounds
on the number $m_\times$ of marked edges between $A$ and $W$ are 48
for Ore-degree at most 13 and 41 for Ore-degree at most 12. For each
of the 51 arrangements of the spokes in
\textbf{Figures~\ref{fig:finite-seven-disjoint},
\ref{fig:finite-seven-separate} and~\ref{fig:finite-seven-shared}},
fix all edges incident with the roots. The encoding allows 16 vertices
in $W$, with unused vertices isolated.

In each formula the variables $E_{xy},M_{xy},J_{xy}$
indicate, respectively, that $x,y$ are adjacent in the host,
that at least one marked edge joins them and that two marked edges
join them. We require
$J_{xy}\Rightarrow M_{xy}\Rightarrow E_{xy}$. The host
and marked degrees at a vertex are sums of $E_{xy}+J_{xy}$
and $M_{xy}+J_{xy}$ over the pairs containing it. For each
pair of disjoint candidate marked edges $uv,xy$, impose
\begin{equation}\label{eq:multi-sat-witness}
 \neg M_{uv}\ \lor\ \neg M_{xy}\ \lor\
 E_{ux}\ \lor\ E_{uy}\ \lor\ E_{vx}\ \lor\ E_{vy}.
\end{equation}
Thus two disjoint marked edges require a host edge between their
endpoint sets. Such a joining edge may be unmarked. The
formulas also impose the degree bounds and the capacity and
Ore-degree inequalities of \textbf{Subsection~\ref{subsec:exact-small-degree}}, using Boolean variables that
indicate whether the relevant counts exceed fixed
thresholds.
For each arrangement of spokes the program permits an ordering of the
remaining vertices whose adjacency data are lexicographically maximal
under permutations of interchangeable vertices. This justifies the
symmetry conditions comparing the assignment with each transposition
of two successive interchangeable labels.

The generator \path{d7_parallel_sat.py} and the replay
script \path{d7_verify_parallel.py} reconstruct the
51 formulas in conjunctive normal form (CNF) from
\path{d7_parallel_sat_cuts_q9} and
\path{d7_parallel_sat_cuts_q10}. The replay compares the
reconstructed formulas to the archived formulas and uses
\path{d7_rup_replay.c} to check 22,217,393 clause additions.
The checker uses reverse unit propagation. For each added
clause, it assumes the negation of each literal and uses
the indicated clauses still present to obtain a
contradiction by unit propagation. Each replay derives
the empty clause, excluding the corresponding spoke
configuration.

\subsubsection{Maximum Degree Nine}\label{app:multi-nine}

The repository~\cite{FGSTCode} contains the degree-nine reductions
used in \textbf{Lemma~\ref{lem:finite-nine}}. Its index is
\path{reproducibility/multigraph/REPLAY.md}. For the step $q\geq13$, see
Sections~1, 2 and 4--12 of \path{delta9_parallel_reduction.md}.
Sections~12 and 14--17 give the bounds on the components of the
graph of double pairs and the restrictions on their incident edges.
The entry \emph{Degree nine, multiplicities and repeated roots} in
the index lists the separate argument for the three component
counts when $b=18$. Section~18 of the initial reduction excludes
$n=22$ in the marked-triangle case.
The index also gives the equality check for 97 distinct marked
pairs, the five-cycle inequality when $b=1$ and there is no marked
triangle, and the reductions for a host of order 23. Sections~2
and~3 of the initial reduction give the marked-triangle degree counts.
The proofs linked from the entries for slack seven and nine also
give the lower bound on the number of endpoints of unmarked edges
when the marked graph is triangle-free.
For $q=14$ with all spokes marked, we use the notation $X,Y,J,B_0,U$
of \textbf{Lemma~\ref{lem:finite-nine}}. If $B_0$ is connected,
\path{d9_q14_connected_spokes.py} tests the 102 pairs $(U,J)$ in
\textbf{Figures~\ref{fig:nine-qfourteen-connected}
and~\ref{fig:nine-qfourteen-connected-last}}. An outside vertex with
marked neighbor set $S\subseteq A$ needs at least
$|N_{B_0}(S)\setminus S|$ unmarked neighbors in $A$. Each $a\in A$
can contribute at most $8-d_J(a)$ further incidences to the outside
vertices. The integer recurrence first imposes the total budget $U$,
then the capacity of one and of two vertices of $A$, excluding
2, 64 and 11 instances, respectively. For the remaining 25 instances
it enumerates the host neighbor set containing $S\cup N_{B_0}(S)$
and the possible extra marked copies for each outside vertex.
Each such vertex uses at least one unit of $U$.
Summing over the possible multisets of outside vertices excludes
all 25 instances.

If $B_0$ is disconnected, the verifier is
\path{d9_q14_split_spokes.py}. When a component has order at least
ten, the routine \texttt{large\_component} first restricts the two
part sizes and the numbers of marked and unmarked edges to that
component. The only remaining pair of part sizes is $(5,5)$.
Section~16 of \path{delta9_parallel_reduction.md} gives the final
degree count. One vertex would need degree nine but has only seven.
When all components have order at most nine, the 33 profiles in
\textbf{Figure~\ref{fig:nine-qfourteen-split}} give 58 choices with
ordered parts and 174 after specifying $U$. For each choice the
routine enumerates the marked and unmarked neighbors that an outside
vertex can have in each part of each component. It sums the possible
contributions subject to the degree bounds on both parts and the
total budget $U$. In all 174 instances the resulting upper bound
is smaller than $52+U$.

For the order-23 cases with a marked triangle, the generator
\path{d9_blue_ferrers_sat.py} builds 112 Boolean formulas from
the configurations in \textbf{Figures~\ref{fig:nine-chain-1}--\ref{fig:nine-chain-6}}.
Each formula uses the nested rows and columns of the chain matrix,
the degree bounds and \eqref{eq:multi-sat-witness}. The instances
are stored in \path{chain_n23_broad}.
The LRAT refutations of the 112 formulas contain 8,183,267 clause
additions in total. An LRAT file lists clauses and
gives hints for checking them by unit propagation. The 18 instances
with no marked triangle are generated by \path{d9_order23_blue_chain_sat.py}.
Each occurs at a starred panel among the same 112 templates. Their
refutations have 3,657,159 clause additions. The auxiliary
checks \path{n23_t2_same} and \path{n23_t2_distinct} treat the two ways
in which two exterior vertices can each have two neighbors on the
marked triangle.

When $b=1$ and there is no marked triangle, 33 of the 60 possible
slack vectors with $R=1,3,5,7$ are excluded by the capacity bounds.
The 27 other cases are covered by \path{d9_anchor_cases.py} and
\path{d9_b1_anchor_proofs}, with 1,726,943 clause additions in total.
When there is a marked triangle the remaining vector is $(1,1,2,1,2)$.
The program \path{d9_pentagon.py} checks the refutation
\path{d9_b1_triangle/s11212}, with 211,080 clause additions.

For the minimum-slack test, rotate and reflect each vector
$(s_0,\ldots,s_4)$ and keep one copy of each orbit for which all five
part sizes are positive. There are 38 vectors for $R=7$. The bounds
at the cycle edges exclude 19 and the bounds with all $s_i>0$
exclude another three. The corresponding counts for $R=9$ are 54
initial vectors and ten exclusions. Each remaining vector gives two
instances, according as the marked edges contain a triangle or not.
The formulas impose the anchor signatures, the adjacencies between
successive parts, the mark and degree constraints, the range of
$b$ and the conditions for minimum slack in \textbf{Lemma~\ref{lem:finite-nine}}.

The generator \path{d9_pentagon.py} reconstructs the CNF
for each slack vector and choice of whether a marked
triangle exists. For $R=7$, it checks 32 refutations with
1,343,507 clause additions. For $R=9$, it checks 88
refutations with 17,817,224 clause additions. The fixed
directories have prefixes \path{d9_r7_zero_triangle}
and \path{d9_r9_triangle}. The replay compares the
uncompressed CNF and LRAT hashes with the archived
values.
The resulting empty clauses exclude every remaining instance at $R=7,9$.
The file \path{reproducibility/multigraph/REPLAY.md} indexes the
generators, parameters, fixed instances and refutations.
